% Sustitución provisional del bloque de Fase 5, hasta la prueba con banda detenida.
% No modificar las Fases 1--4 ni presentar esta entrega como el cierre de la Fase 5.

En la quinta fase se inició la integración de la percepción visual y del posicionamiento automático del brazo. Antes de incorporar la medición del movimiento de la banda mediante el codificador rotativo, se desarrollaron cuatro subetapas: preparación del conjunto de datos y entrenamiento del modelo, integración de la comunicación con la cámara, calibración geométrica del plano de trabajo y comprobación del posicionamiento con la banda detenida.

Para preparar el conjunto de datos, se capturaron fotografías y videos de las piezas en distintas posiciones. De los videos se seleccionaron fotogramas, que se incorporaron junto con las fotografías al material de entrenamiento. Se utilizó Mind+ para cargar las imágenes y preparar las anotaciones en formato YOLO. Se delimitaron manualmente las piezas mediante cajas delimitadoras y se reunieron 400 imágenes por tipo de pieza, para un total de 800 imágenes. El entrenamiento se orientó a detectar las piezas, localizar sus regiones en la imagen y diferenciar los tipos 6 y 7. El formato de las anotaciones se distinguió de la arquitectura del modelo, cuya identificación quedó pendiente en los registros de entrenamiento.

Se examinó cualitativamente la detección y la asignación de clases. A partir de las dificultades observadas en la diferenciación entre piezas, se delimitó el alcance de la aplicación a la localización para la manipulación, sin considerar validada la clasificación por tipo. Según la descripción del autor, se llevó un conteo global de piezas recogidas, separado de la clasificación; su procedimiento, ubicación dentro de la secuencia de pruebas y resultados numéricos quedaron pendientes de documentación.

La comunicación entre la cámara y la ESP32 se probó inicialmente mediante I2C. Durante la integración se observaron, según el registro del autor, incidencias en la actualización de la OLED o en el intercambio con la Portenta. Posteriormente, el enlace de la cámara se trasladó a UART2 y se verificó su interacción mediante mensajes de terminal. Se mantuvieron los buses I2C de la pantalla y del enlace entre microcontroladores descritos en la cuarta fase. Las consultas a la biblioteca de la cámara se concentraron en una tarea de FreeRTOS y se implementó una máquina de estados para conexión, calibración, apertura del modelo y detección.

Para establecer la correspondencia entre las coordenadas de imagen y el plano de trabajo, se utilizaron cuatro marcadores fiduciales AprilTag colocados en posiciones conocidas. Su uso se correspondió con la referencia visual mediante identificadores, documentada para este tipo de marcadores \cite{apriltag_fase5}. Se asignaron coordenadas métricas a sus centros a partir del ancho útil de banda de 292~mm, el ancho total de 412~mm y la separación longitudinal de 382~mm entre filas de marcadores, confirmados por el autor. La cámara se colocó en el modo de reconocimiento de etiquetas y se acumularon 25 observaciones del centro de cada marcador. Se promediaron las coordenadas de imagen y se utilizaron las cuatro correspondencias para resolver una homografía plana.

La transformación se calculó mediante un sistema lineal de ocho incógnitas y se aplicó a las coordenadas de imagen de las detecciones. Se comprobaron la disponibilidad de las muestras, la validez numérica de la solución y la pertenencia de las posiciones al área útil configurada. El procedimiento se trató como una calibración geométrica del plano de trabajo; no incluyó una estimación de los parámetros intrínsecos ni de la distorsión óptica de la cámara. Después del cálculo se solicitó la apertura del modelo personalizado y se habilitó la lectura de piezas.

Con la banda detenida, se realizaron cinco pruebas de posicionamiento. En cada prueba se colocó una pieza y se comprobó que el brazo se desplazara hasta la ubicación correspondiente, sin ejecutar su recogida. En el software se filtraron las detecciones, se transformaron las posiciones al sistema de referencia del brazo y se enviaron destinos a las funciones de posicionamiento X/Y. La aceptación se basó en la coincidencia de ubicación reportada por el autor; no se documentaron coordenadas de comparación, un instrumento de medición ni una tolerancia numérica. Esta comprobación funcional precedió a la incorporación del encoder y se distinguió de los ensayos posteriores con la banda en movimiento.

% PENDIENTE: arquitectura, parámetros y particiones de entrenamiento; fechas;
% fotografía de los tags y comprobación métrica del origen cámara/brazo.
% La integración del encoder y la sincronización dinámica se redactarán después.
